% Generated draft: requires mathematical and source review.
\documentclass[11pt]{article}
\usepackage[margin=25mm]{geometry}
\usepackage{fontspec}
\setmainfont{Latin Modern Roman}
\usepackage{amsmath,amssymb,mathtools,unicode-math}
\setmathfont{Latin Modern Math}
\usepackage{enumitem,needspace,xurl,hyperref}
\setlength{\parindent}{0pt}
\setlength{\parskip}{5pt}
\setlength{\emergencystretch}{3em}
\newcommand{\N}{\mathbb{N}}
\newcommand{\Z}{\mathbb{Z}}
\newcommand{\Q}{\mathbb{Q}}
\newcommand{\R}{\mathbb{R}}
\newcommand{\C}{\mathbb{C}}
\newcommand{\F}{\mathbb{F}}
\newcommand{\A}{\mathbb{A}}
\newcommand{\PP}{\mathbb{P}}
\newcommand{\E}{\mathbb{E}}
\newcommand{\eps}{\varepsilon}
\newcommand{\dd}{\,\mathrm{d}}
\begin{document}
\section*{UnsolvedMath ID: 78}
\textbf{Source problem code:} catalogue entry
\subsection{Definitions and mathematical statement}
DARPA's Mathematical Challenge Nine (2007), ``Physical Consequences of Perelman's Proof,'' asks for applications of Perelman's proof of the Poincar\'e conjecture to materials fabrication across scales, from nanostructures to macroscopic fabrication. The background frames this as asking how the breakthrough in three-dimensional geometry can inform materials science and refers more broadly to consequences of geometrization.

The requested task is to identify and formulate a mechanism connecting a geometrization invariant or other explicitly specified consequence of the proof to a fabrication design or physical prediction, together with a concrete materials objective and a way to test that mechanism. The source does not select such an invariant, a material or geometric model, a mechanism, scales, experimental conditions, or a quantitative objective. This description therefore records an application agenda rather than a quantified mathematical proposition.

The local registry marks the record open. The supplied literature triage dated 2026-08-17 classifies it as partial progress, states that Perelman's proof is established, and says that a solution or refutation of this application agenda is not applicable without a concrete materials objective. It does not assert that the Poincar\'e conjecture itself remains open.
\subsection{Short English statement}
Find a concrete, testable way to use Perelman's work on three-dimensional geometry in materials fabrication from nano- to macroscopic scales. The supplied agenda lacks a specified physical mechanism and materials objective, so it does not yet define a theorem to prove.
\subsection{Sources}
\begin{itemize}
\item DARPA DSO, Broad Agency Announcement 07-68, Mathematical Challenge Nine (2007), supplied primary reference: \url{https://web.math.utk.edu/~vasili/refs/darpa07.MathChallenges.html}.
\item Local registry data/unsolvedmath/problems.json, permanent catalogue entry, source code catalogue entry, ``Physical Consequences of Perelman's Proof,'' attributed to DARPA (2007); complete statement, background, and literature triage dated 2026-08-17. Dataset attribution: Ulam AI and the UnsolvedMath Contributors.
\end{itemize}
\end{document}
